\section*{Problems}

\subsection*{Problem Statements}

% Remember
\begin{problema}
\label{prob:0d6ff67}
Urban-car fuel efficiency ($Y$, km/l) is studied as a function of two continuous predictors: gross vehicle weight ($X_1$, tons) and engine displacement ($X_2$, liters). From a random sample of $n=30$ cars under standardized conditions, the OLS multiple-regression model is
\begin{align}
\hat Y=25.0-3.0X_1-1.5X_2,\quad R^2=0.780,\quad R_{\text{adj}}^2=0.763.
\end{align}
\begin{enumerate}
\item Interpret the partial coefficients $\hat\beta_1=-3.0$ and $\hat\beta_2=-1.5$, explicitly stating the ceteris-paribus condition.
\item Calculate the point prediction for a compact sedan weighing 1.5 tons with a 2.0-liter engine.
\item If the total sample variance of observed efficiencies is $S_Y^2=12.0$, calculate the Regression Sum of Squares ($\text{SSR}$) and Residual Sum of Squares ($\text{SSE}$).
\end{enumerate}
\end{problema}

% Understand
\begin{problema}
\label{prob:9f8863e}
A marketing-consulting firm analyzes the impact of advertising-channel diversification on monthly corporate sales ($Y$, thousands of dollars). In a historical sample of $n=40$ quarters fitted with three independent variables ($p=3$), explained variance is $R^2=0.780$.
\begin{enumerate}
\item Calculate adjusted $R^2=1-\left[\dfrac{n-1}{n-p-1}(1-R^2)\right]$.
\item Explain why, in multivariable analysis, adjusted $R^2$ is strictly less than nominal $R^2$ whenever two or more predictors are included.
\item If an analyst adds an irrelevant fourth predictor $X_4$ (headquarters ZIP code), with $R^2$ barely increasing from $0.780$ to $0.781$, determine the new adjusted $R^2$ for $p=4$.
\end{enumerate}
\end{problema}

% Apply
\begin{problema}
\label{prob:1d085e7}
A multiple-regression model predicts residential appraisal value ($Y$, thousands of dollars) from $p=4$ variables: built area ($X_1$, m$^2$), number of rooms ($X_2$), age ($X_3$, years), and distance from the city center ($X_4$, km). For $n=50$ properties,
\begin{align}
\hat Y=50.0+0.80X_1+15.0X_2-2.0X_3-5.0X_4,
\end{align}
with $\text{VIF}(X_1)=4.2$, $\text{VIF}(X_2)=3.8$, $\text{VIF}(X_3)=1.5$, and $\text{VIF}(X_4)=1.2$.
\begin{enumerate}
\item Interpret the economic meaning of the area ($\hat\beta_1$) and distance ($\hat\beta_4$) coefficients.
\item Evaluate collinearity diagnostics using the VIFs: is there severe multicollinearity or a moderately admissible level?
\item Calculate the projected appraisal value for a residence with 120 m$^2$, 3 rooms, 10 years of age, and 8 km from the center.
\item What modeling strategy would you apply if $\text{VIF}(X_1)$ exceeded the critical threshold of 10?
\end{enumerate}
\end{problema}

% Analyze
\begin{problema}
\label{prob:15ced94}
\textbf{Analytical derivation of vector-matrix OLS and population unbiasedness in the general $p$-regressor model:} consider $\mathbf Y=\mathbf X\boldsymbol\beta+\boldsymbol\varepsilon$, where $\mathbf X$ has full column rank, $E[\boldsymbol\varepsilon]=\mathbf0$, and $\text{Var}(\boldsymbol\varepsilon)=\sigma^2\mathbf I_n$.
\begin{enumerate}
\item Starting from $S(\boldsymbol\beta)=(\mathbf Y-\mathbf X\boldsymbol\beta)^T(\mathbf Y-\mathbf X\boldsymbol\beta)$, use matrix differential calculus to derive the \textbf{normal equations} $\mathbf X^T\mathbf X\hat{\boldsymbol\beta}=\mathbf X^T\mathbf Y$ and justify the uniqueness of $\hat{\boldsymbol\beta}=(\mathbf X^T\mathbf X)^{-1}\mathbf X^T\mathbf Y$.
\item Prove that $\hat{\boldsymbol\varepsilon}=(\mathbf I_n-\mathbf H)\mathbf Y$ is orthogonal to the column space of $\mathbf X$ ($\mathbf X^T\hat{\boldsymbol\varepsilon}=\mathbf0$) and that $\hat{\mathbf Y}^T\hat{\boldsymbol\varepsilon}=0$.
\item Prove that $\hat{\boldsymbol\beta}$ is unbiased and that $\text{Var}(\hat{\boldsymbol\beta})=\sigma^2(\mathbf X^T\mathbf X)^{-1}$.
\end{enumerate}
\end{problema}

% Evaluate
\begin{problema}
\label{prob:ee59291}
A data scientist models annual salary ($Y$, thousands of dollars) from years of experience ($X_1$) and education indicators, with Bachelor's degree as the reference group ($D_1=D_2=0$): $D_1=1$ for a Bachelor's degree and $D_2=1$ for a Master's degree. For $n=100$ employees, $\hat Y=22.0+2.5X_1+12.0D_1+25.0D_2$, with $R^2=0.760$.
\begin{enumerate}
\item Interpret the intercept $\hat\beta_0=22.0$ and categorical coefficients $\hat\beta_2=12.0$ and $\hat\beta_3=25.0$.
\item Write the three salary-versus-experience lines by education group, and predict salary for a manager with 8 years of experience and a Master's degree.
\item Explain the \textbf{Dummy Variable Trap} and why including a third indicator $D_0=1$ for Bachelor's degree together with the intercept would be fatal to the model's matrix algebra.
\end{enumerate}
\end{problema}

% Create
\begin{problema}
\label{prob:a6a21c3}
Design an original multiple-regression problem (different from cars, marketing, housing, or salaries): a hospital models total hospitalization cost ($Y$, thousands of dollars) from patient age ($X_1$, years) and length of stay ($X_2$, days). For $n=45$ patients, $\hat Y=2.0+0.05X_1+1.2X_2$, $R^2=0.71$, $SE(\hat\beta_1)=0.02$, and $SE(\hat\beta_2)=0.15$. Formulate and solve individual $t$ tests for both predictors ($\nu=42$, $t_{0.025,42}=2.018$), interpret the model, and predict cost for a 65-year-old patient with a 5-day stay.
\end{problema}

\subsection*{Detailed Solutions}

% Remember
\begin{solproblema}[prob:0d6ff67]
\begin{enumerate}
\item Holding $X_2$ constant, each additional ton of weight ($X_1$) reduces efficiency by 3.0 km/l. Holding $X_1$ constant, each additional liter of displacement ($X_2$) reduces efficiency by 1.5 km/l.
\item $\hat Y=25.0-3.0(1.5)-1.5(2.0)=17.5$ km/l.
\item $\text{SST}=(n-1)S_Y^2=(29)(12.0)=348.0$. $\text{SSR}=R^2\times\text{SST}=(0.780)(348.0)=271.44$. $\text{SSE}=348.0-271.44=76.56$.
\end{enumerate}
\end{solproblema}

% Understand
\begin{solproblema}[prob:9f8863e]
\begin{enumerate}
\item $R_{\text{adj}}^2=1-\left[\frac{39}{36}(0.220)\right]=1-0.23833\approx0.7617$ (76.17\%).
\item Adding a predictor can never increase $\text{SSE}$ ($\text{SSE}_{p+1}\le\text{SSE}_p$), so nominal $R^2$ never decreases, encouraging overfitting. Adjusted $R^2=1-\text{MSE}/\text{MST}$ corrects this: adding a predictor decreases $(n-p-1)$ and penalizes the MSE/MST ratio unless the new variable explains enough real variance.
\item With $R^2_{\text{new}}=0.781$, $p=4$, $R_{\text{adj,new}}^2=1-\left[\frac{39}{35}(0.219)\right]\approx0.7560$. Although nominal $R^2$ rose slightly, adjusted $R^2$ fell from 0.7617 to 0.7560, showing that ZIP code adds no useful explanatory value.
\end{enumerate}
\end{solproblema}

% Apply
\begin{solproblema}[prob:1d085e7]
\begin{enumerate}
\item $\hat\beta_1=0.80$: controlling for rooms, age, and distance, each additional m$^2$ increases value by \$800. $\hat\beta_4=-5.0$: holding the other variables fixed, each additional kilometer from the center reduces value by \$5,000.
\item Diagnostic scale: VIF below 5 is low/moderate, $5\le\text{VIF}<10$ is concerning, and VIF at least 10 is severe. The largest VIF is $4.2<5$, so there is \textbf{no severe multicollinearity}; the moderate overlap between area and rooms is natural, with enough orthogonal variation to estimate both effects.
\item $\hat Y=50.0+0.80(120)+15.0(3)-2.0(10)-5.0(8)=131.0$ thousand dollars.
\item If VIF exceeded 10, combine collinear variables into a synthetic index, remove the variable with the least significant $t$, use principal-component regression, or use Ridge regression $\hat{\boldsymbol\beta}_{\text{Ridge}}=(\mathbf X^T\mathbf X+\lambda\mathbf I)^{-1}\mathbf X^T\mathbf Y$ to stabilize matrix inversion.
\end{enumerate}
\end{solproblema}

% Analyze
\begin{solproblema}[prob:15ced94]
\begin{enumerate}
\item Expanding $S(\boldsymbol\beta)=\mathbf Y^T\mathbf Y-2\boldsymbol\beta^T\mathbf X^T\mathbf Y+\boldsymbol\beta^T\mathbf X^T\mathbf X\boldsymbol\beta$ and differentiating,
\begin{align}
\frac{\partial S(\boldsymbol\beta)}{\partial\boldsymbol\beta}=-2\mathbf X^T\mathbf Y+2\mathbf X^T\mathbf X\boldsymbol\beta=\mathbf0\implies\mathbf X^T\mathbf X\hat{\boldsymbol\beta}=\mathbf X^T\mathbf Y.
\end{align}
Since $\mathbf X$ has full rank, $\mathbf X^T\mathbf X$ is positive definite and nonsingular, guaranteeing the unique solution $\hat{\boldsymbol\beta}=(\mathbf X^T\mathbf X)^{-1}\mathbf X^T\mathbf Y$.
\item With $\hat{\boldsymbol\varepsilon}=\mathbf Y-\mathbf X\hat{\boldsymbol\beta}$, $\mathbf X^T\hat{\boldsymbol\varepsilon}=\mathbf X^T\mathbf Y-\mathbf X^T\mathbf X\hat{\boldsymbol\beta}=\mathbf0$ by the normal equations. Since $\hat{\mathbf Y}=\mathbf X\hat{\boldsymbol\beta}$ lies in the column space of $\mathbf X$, $\hat{\mathbf Y}^T\hat{\boldsymbol\varepsilon}=\hat{\boldsymbol\beta}^T\mathbf X^T\hat{\boldsymbol\varepsilon}=0$.
\item Substituting $\mathbf Y=\mathbf X\boldsymbol\beta+\boldsymbol\varepsilon$ gives $\hat{\boldsymbol\beta}=\boldsymbol\beta+(\mathbf X^T\mathbf X)^{-1}\mathbf X^T\boldsymbol\varepsilon$. Taking expectations gives $E[\hat{\boldsymbol\beta}]=\boldsymbol\beta$. For the variance,
\begin{align}
\text{Var}(\hat{\boldsymbol\beta})=(\mathbf X^T\mathbf X)^{-1}\mathbf X^T E[\boldsymbol\varepsilon\boldsymbol\varepsilon^T]\mathbf X(\mathbf X^T\mathbf X)^{-1}=\sigma^2(\mathbf X^T\mathbf X)^{-1}.
\end{align}
\end{enumerate}
\end{solproblema}

% Evaluate
\begin{solproblema}[prob:ee59291]
\begin{enumerate}
\item $\hat\beta_0=22.0$: estimated salary of the Bachelor's reference group at $X_1=0$. $\hat\beta_2=12.0$: at equal experience, a Bachelor's graduate earns \$12,000 more than the reference group. $\hat\beta_3=25.0$: at equal experience, a Master's graduate earns \$25,000 more than the reference group.
\item Bachelor's: $\hat Y=22.0+2.5X_1$. Licentiate degree: $\hat Y=34.0+2.5X_1$. Master's: $\hat Y=47.0+2.5X_1$. For $X_1=8$ and a Master's degree, $\hat Y=67.0$ thousand dollars.
\item Including all three indicators ($D_0,D_1,D_2$) with an intercept gives $D_0+D_1+D_2=\mathbf1_n$ in every row, making the columns of $\mathbf X$ linearly dependent and $\det(\mathbf X^T\mathbf X)=0$. The normal matrix is singular, so $(\mathbf X^T\mathbf X)^{-1}\mathbf X^T\mathbf Y$ cannot be computed; one category must be omitted as the baseline.
\end{enumerate}
\end{solproblema}

% Create
\begin{solproblema}[prob:a6a21c3]
$H_{0,1}:\beta_1=0$ versus $H_{a,1}:\beta_1\ne0$; $H_{0,2}:\beta_2=0$ versus $H_{a,2}:\beta_2\ne0$. The statistics are
\begin{align}
t_1=\frac{0.05}{0.02}=2.5,\qquad t_2=\frac{1.2}{0.15}=8.0.
\end{align}
Since $|t_1|=2.5>2.018$ and $|t_2|=8.0>2.018$, \textbf{reject both} null hypotheses: age and length of stay each contribute significantly to cost while controlling for the other. Holding length of stay fixed, each additional year of age increases cost by \$50; holding age fixed, each additional day increases cost by \$1,200. For a 65-year-old with a 5-day stay,
\begin{align}
\hat Y=2.0+0.05(65)+1.2(5)=11.25\text{ thousand dollars.}
\end{align}
\end{solproblema}
